%---------------------------------------------------------------------
% Urdu edition, units 1-4.
%---------------------------------------------------------------------

\section*{یونٹ \textenglish{1}: احتمال}
\markboth{یونٹ 1: احتمال}{}

\begin{nukta}[چار اصول جن پر پورا یونٹ کھڑا ہے]
\begin{ltrtab}
\begin{minipage}{\linewidth}
\begin{align*}
\text{متمم}          &: && \prob{A'}=1-\prob{A}\\
\text{جمع (عام)}     &: && \prob{A\cup B}=\prob{A}+\prob{B}-\prob{A\cap B}\\
\text{ضرب (آزاد)}    &: && \prob{A\cap B}=\prob{A}\,\prob{B}\\
\text{مشروط}         &: && \prob{A\mid B}=\frac{\prob{A\cap B}}{\prob{B}}
\end{align*}
\end{minipage}
\end{ltrtab}
اگر دو واقعات \emph{باہم خارج} ہوں تو \(\prob{A\cap B}=0\) اور جمع کا اصول
سادہ ہو کر \(\prob{A}+\prob{B}\) رہ جاتا ہے۔
\end{nukta}

\begin{sawal}{سوال \textenglish{1}}
پائلٹ اسکول میں داخلے کا احتمال \(0.5\) ہے۔ اگر تین امیدوار بے ترتیب چنے جائیں
تو احتمال نکالیے کہ
\textbf{(الف)} تینوں کو داخلہ ملے،
\textbf{(ب)} کسی کو نہ ملے،
\textbf{(ج)} صرف ایک کو ملے۔
\end{sawal}

\begin{hal}
تینوں امیدوار ایک دوسرے سے آزاد ہیں، اس لیے احتمالات آپس میں ضرب ہوں گے۔

\medskip
\textbf{(الف)} تینوں کو داخلہ ملنے کی صرف ایک صورت ہے:
\[
(0.5)(0.5)(0.5)=(0.5)^{3}=0.125
\]

\textbf{(ب)} کسی کو نہ ملنے کی بھی ایک ہی صورت ہے:
\[
(0.5)^{3}=0.125
\]

\textbf{(ج)} ''صرف ایک'' میں یہ بھی گننا ہے کہ وہ ایک \emph{کون} ہے --- پہلا،
دوسرا یا تیسرا۔ یعنی تین صورتیں:
\[
\binom{3}{1}(0.5)^{1}(0.5)^{2}=3(0.125)=0.375
\]

\parakh{چاروں صورتیں \(0,1,2,3\) مل کر پوری نمونہ فضا بنتی ہیں:
\(0.125+0.375+0.375+0.125=1\)۔}

\jawab{(الف) \(0.125\) \quad (ب) \(0.125\) \quad (ج) \(0.375\)}
\end{hal}

\begin{ghalti}
حصہ (ج) میں \(\binom{3}{1}=3\) سے ضرب لازمی ہے۔ (الف) اور (ب) میں یہ ضرب نہیں
لگتی، کیونکہ ''سب'' یا ''کوئی نہیں'' کی صرف ایک ہی ترتیب ممکن ہے۔
\end{ghalti}

\begin{sawal}{سوال \textenglish{2}}
ایک مرتبان میں \textenglish{8} سبز نقطہ دار، \textenglish{10} سبز دھاری دار،
\textenglish{12} نیلی نقطہ دار اور \textenglish{10} نیلی دھاری دار گیندیں ہیں۔
ایک گیند بے ترتیب نکالی جائے تو احتمال کیا ہے کہ وہ
\textbf{(الف)} سبز یا دھاری دار ہو، \textbf{(ب)} نقطہ دار ہو،
\textbf{(ج)} نیلی یا نقطہ دار ہو؟
\end{sawal}

\begin{hal}
سب سے پہلے اعداد کو دو رخی جدول میں رکھیے۔ اس سوال کے تقریباً سارے نمبر اسی قدم
پر ملتے یا کھوتے ہیں۔

\medskip
\begin{center}
\begin{tabular}{@{}p{2.6cm} p{2.4cm} p{2.4cm} p{2.0cm}@{}}
\toprule
 & \textbf{نقطہ دار} & \textbf{دھاری دار} & \textbf{کل} \\
\midrule
سبز  & \textenglish{8}  & \textenglish{10} & \textbf{\textenglish{18}} \\
نیلی & \textenglish{12} & \textenglish{10} & \textbf{\textenglish{22}} \\
\midrule
\textbf{کل} & \textbf{\textenglish{20}} & \textbf{\textenglish{20}} & \textbf{\textenglish{40}} \\
\bottomrule
\end{tabular}
\end{center}

\medskip
\textbf{(الف)} سبز اور دھاری دار میں \emph{اشتراک} ہے (\textenglish{10} گیندیں
دونوں میں ہیں)، اس لیے عام جمع کا اصول لگے گا:
\[
\prob{G\cup S}=\frac{18}{40}+\frac{20}{40}-\frac{10}{40}=\frac{28}{40}=0.70
\]

\textbf{(ب)} سیدھا ستون کے کل سے:
\[
\prob{D}=\frac{20}{40}=0.50
\]

\textbf{(ج)} اشتراک \textenglish{12} نیلی نقطہ دار گیندیں ہیں:
\[
\prob{B\cup D}=\frac{22}{40}+\frac{20}{40}-\frac{12}{40}=\frac{30}{40}=0.75
\]

\parakh{براہِ راست گنتی سے: سبز یا دھاری دار \(=8+10+10=28\)؛ نیلی یا نقطہ دار
\(=12+10+8=30\)۔ دونوں مطابقت رکھتے ہیں۔}

\jawab{(الف) \(0.70\) \quad (ب) \(0.50\) \quad (ج) \(0.75\)}
\end{hal}

\begin{sawal}{سوال \textenglish{3}}
اگر \(\prob{E}=0.2\)، \(\prob{F}=0.5\) اور \(\prob{E\cup F}=0.6\) ہو تو نکالیے
\textbf{(الف)} \(\prob{E\cap F}\)، \textbf{(ب)} \(\prob{E\mid F}\)،
\textbf{(ج)} \(\prob{F\mid E}\)۔
\end{sawal}

\begin{hal}
\textbf{(الف)} جمع کے عام اصول کو الٹ کر اشتراک کو مطلوب بنائیے:
\[
\prob{E\cap F}=\prob{E}+\prob{F}-\prob{E\cup F}=0.2+0.5-0.6=0.1
\]

\textbf{(ب)} \(\displaystyle \prob{E\mid F}=\frac{0.1}{0.5}=0.2\)

\textbf{(ج)} \(\displaystyle \prob{F\mid E}=\frac{0.1}{0.2}=0.5\)

\medskip
\textbf{ایک اضافی نکتہ جس کے نمبر ملتے ہیں۔}
غور کیجیے کہ \(\prob{E\mid F}=0.2=\prob{E}\) اور \(\prob{F\mid E}=0.5=\prob{F}\)
ہے۔ یعنی ایک واقعے کا علم دوسرے کے بارے میں کچھ نہیں بتاتا۔ برابر کی بات:
\(\prob{E}\prob{F}=(0.2)(0.5)=0.1=\prob{E\cap F}\)۔ لہٰذا \(E\) اور \(F\)
\textbf{آزاد} واقعات ہیں۔

\jawab{(الف) \(0.1\) \quad (ب) \(0.2\) \quad (ج) \(0.5\)؛ اور \(E,F\) آزاد ہیں۔}
\end{hal}

%---------------------------------------------------------------------
\section*{یونٹ \textenglish{2}: بے ترتیب متغیر}
\markboth{یونٹ 2: بے ترتیب متغیر}{}

\begin{nukta}[کوئی فہرست احتمالی تقسیم کب کہلاتی ہے]
\[
0\le P(X=x)\le 1 \quad\text{ہر } x \text{ کے لیے},
\qquad
\sum_{x} P(X=x)=1
\]
دوسری شرط ہی مفقود قدر والے سوالات کا پورا حل ہے، اور باقی سوالات میں حساب کی
غلطی پکڑنے کا مفت ذریعہ۔
\end{nukta}

\begin{sawal}{سوال \textenglish{4}}
ایک کھرا سکہ تین بار اچھالا جاتا ہے۔ \(X\) سے مراد چِتوں کی تعداد ہے۔ احتمالی
تقسیم بنائیے اور \(P(X\ge 2)\) نکالیے۔
\end{sawal}

\begin{hal}
کل \(2^{3}=8\) یکساں ممکن نتائج ہیں۔ انہیں لکھ دینے سے گنتی مشکوک نہیں رہتی:

\medskip
\begin{center}
\begin{tabular}{@{}p{2.0cm} p{5.4cm} p{1.8cm} p{2.0cm}@{}}
\toprule
\(x\) & نتائج & تعداد & \(P(X=x)\) \\
\midrule
\textenglish{0} & \textenglish{TTT}                & \textenglish{1} & \(1/8\) \\
\textenglish{1} & \textenglish{HTT, THT, TTH}      & \textenglish{3} & \(3/8\) \\
\textenglish{2} & \textenglish{HHT, HTH, THH}      & \textenglish{3} & \(3/8\) \\
\textenglish{3} & \textenglish{HHH}                & \textenglish{1} & \(1/8\) \\
\midrule
\textbf{کل} & & \textbf{\textenglish{8}} & \textbf{\textenglish{1}} \\
\bottomrule
\end{tabular}
\end{center}

\medskip
''دو یا زیادہ'' کا مطلب \(X=2\) یا \(X=3\) ہے، اور یہ باہم خارج ہیں:
\[
P(X\ge 2)=\frac{3}{8}+\frac{1}{8}=\frac{4}{8}=\frac12=0.5
\]

\jawab{تقسیم: \(\tfrac18,\ \tfrac38,\ \tfrac38,\ \tfrac18\)؛ اور
\(P(X\ge2)=0.5\)}
\end{hal}

\begin{sawal}{سوال \textenglish{5}}
ایک تھیلے میں \textenglish{4} سرخ اور \textenglish{6} کالی گیندیں ہیں۔ بغیر
واپسی کے \textenglish{4} گیندوں کا نمونہ لیا جاتا ہے۔ \(X\) سرخ گیندوں کی تعداد
ہو تو \(X\) کی احتمالی تقسیم نکالیے۔
\end{sawal}

\begin{hal}
\emph{بغیر واپسی} کے نمونہ لینے سے یہ دو ذی حدی نہیں بلکہ
\textenglish{hypergeometric} تقسیم بنتی ہے:
\[
P(X=x)=\frac{\dbinom{4}{x}\dbinom{6}{4-x}}{\dbinom{10}{4}},
\qquad \dbinom{10}{4}=210
\]

\medskip
\begin{center}
\begin{ltrtab}
\begin{tabular}{@{}c l c c@{}}
\toprule
$x$ & شمار کنندہ & $P(X=x)$ & اعشاری \\
\midrule
0 & $\binom{4}{0}\binom{6}{4}=15$ & $15/210=1/14$ & 0.0714 \\
1 & $\binom{4}{1}\binom{6}{3}=80$ & $80/210=8/21$ & 0.3810 \\
2 & $\binom{4}{2}\binom{6}{2}=90$ & $90/210=3/7$  & 0.4286 \\
3 & $\binom{4}{3}\binom{6}{1}=24$ & $24/210=4/35$ & 0.1143 \\
4 & $\binom{4}{4}\binom{6}{0}=1$  & $1/210$       & 0.0048 \\
\midrule
\multicolumn{2}{@{}l}{\textbf{کل}} & \textbf{210/210} & \textbf{1.0000} \\
\bottomrule
\end{tabular}
\end{ltrtab}
\end{center}

\parakh{\(15+80+90+24+1=210=\binom{10}{4}\)۔ مجموعہ \(1\) ہے، تقسیم درست ہے۔}

\jawab{\(P(0)=\tfrac1{14}\)، \(P(1)=\tfrac8{21}\)، \(P(2)=\tfrac37\)،
\(P(3)=\tfrac4{35}\)، \(P(4)=\tfrac1{210}\)}
\end{hal}

\begin{ghalti}
اسے دو ذی حدی (\textenglish{binomial}) مان کر \(p=4/10\) لگانے سے
\(P(2)=0.3456\) آتا ہے جو غلط ہے۔ دو ذی حدی کے لیے \(p\) کا ہر بار ایک جیسا
رہنا ضروری ہے، جو صرف \emph{واپسی کے ساتھ} ہوتا ہے۔
\end{ghalti}

\begin{sawal}{سوال \textenglish{6}}
\(X=0,2,4,6,8\) کے احتمالات بالترتیب \(K,\ 2K,\ 4K,\ 2K,\ K\) ہیں۔ \(K\) نکال
کر تقسیم دوبارہ لکھیے۔
\end{sawal}

\begin{hal}
احتمالات کا مجموعہ \(1\) ہونا چاہیے:
\[
K+2K+4K+2K+K=1 \;\Longrightarrow\; 10K=1 \;\Longrightarrow\; K=0.1
\]

\medskip
\begin{center}
\begin{ltrtab}
\begin{tabular}{@{}l c c c c c@{}}
\toprule
$x$      & 0   & 2   & 4   & 6   & 8   \\
$P(X=x)$ & 0.1 & 0.2 & 0.4 & 0.2 & 0.1 \\
\bottomrule
\end{tabular}
\end{ltrtab}
\end{center}

\parakh{\(0.1+0.2+0.4+0.2+0.1=1.0\)، اور ہر قدر \([0,1]\) میں ہے۔ تقسیم
\(x=4\) کے گرد متناسب ہے، اس لیے \(E(X)=4\)۔}

\jawab{\(K=0.1\)؛ تقسیم \(0.1,\ 0.2,\ 0.4,\ 0.2,\ 0.1\)}
\end{hal}

%---------------------------------------------------------------------
\section*{یونٹ \textenglish{3}: مساوات اور عدم مساوات}
\markboth{یونٹ 3: مساوات اور عدم مساوات}{}

\begin{sawal}{سوال \textenglish{7}}
حل کیجیے: \textbf{(الف)} \(t^{2}+4t=21\) (تحلیل سے)، \quad
\textbf{(ب)} \(4x^{2}+3x-1=0\) (دو درجی فارمولے سے)۔
\end{sawal}

\begin{hal}
\textbf{(الف)} پہلے سب کچھ ایک طرف کیجیے --- تحلیل صرف صفر کے مقابل چلتی ہے:
\[
t^{2}+4t-21=0
\]
دو عدد جن کا حاصل ضرب \(-21\) اور جمع \(+4\) ہو: \(+7\) اور \(-3\)۔
\[
(t+7)(t-3)=0 \;\Longrightarrow\; t=-7 \ \text{ یا } \ t=3
\]
\parakh{\((-7)^{2}+4(-7)=49-28=21\)؛ \(3^{2}+4(3)=9+12=21\)۔}

\medskip
\textbf{(ب)} \(a=4,\ b=3,\ c=-1\)۔ پہلے ممیز نکالیے، وہ بتا دیتا ہے کہ جواب کس
قسم کا آئے گا:
\[
b^{2}-4ac=9+16=25>0 \quad\text{(دو حقیقی جڑیں)}
\]
\[
x=\frac{-3\pm\sqrt{25}}{8}=\frac{-3\pm 5}{8}
\;\Longrightarrow\; x=\frac14 \ \text{ یا } \ x=-1
\]

\jawab{(الف) \(t=-7,\ 3\) \quad (ب) \(x=\tfrac14,\ -1\)}
\end{hal}

\begin{sawal}{سوال \textenglish{8}}
دو درجی عدم مساوات حل کیجیے: \(6t^{2}+t-12>0\)
\end{sawal}

\begin{hal}
پہلے جڑیں نکالیے۔ ممیز \(=1+288=289\)، \(\sqrt{289}=17\):
\[
t=\frac{-1\pm 17}{12} \;\Longrightarrow\;
t=\frac{16}{12}=\frac43 \quad\text{یا}\quad t=\frac{-18}{12}=-\frac32
\]
\(t^{2}\) کا عامل \(6>0\) ہے، اس لیے قطع مکافی \emph{اوپر} کو کھلتا ہے: وہ
جڑوں کے \emph{درمیان} محور سے نیچے اور جڑوں سے \emph{باہر} محور سے اوپر ہوتا
ہے۔ ہمیں \(>0\) یعنی اوپر والا حصہ چاہیے:
\[
t<-\frac32 \quad\text{یا}\quad t>\frac43
\]
\parakh{\(t=0\) (جڑوں کے بیچ) رکھیے: \(-12\) آتا ہے، جو \(>0\) نہیں --- درست
طور پر خارج۔ \(t=2\) رکھیے: \(24+2-12=14>0\) --- درست طور پر شامل۔}

\jawab{\(t<-\tfrac32\) یا \(t>\tfrac43\)}
\end{hal}

\begin{sawal}{سوال \textenglish{9}}
حل کیجیے: \textbf{(الف)} \(|2y+5|=|y-4|\)، \quad
\textbf{(ب)} \(|z^{2}-8|\le 8\)
\end{sawal}

\begin{hal}
\textbf{(الف)} \(|u|=|v|\) ہمیشہ دو صورتوں میں بٹتا ہے:

\emph{پہلی صورت:} \(2y+5=y-4 \Rightarrow y=-9\)

\emph{دوسری صورت:} \(2y+5=-(y-4)=-y+4 \Rightarrow 3y=-1 \Rightarrow
y=-\tfrac13\)

\parakh{\(y=-9\): \(|-13|=13\) اور \(|-13|=13\)۔ \(y=-\tfrac13\):
دونوں طرف \(\tfrac{13}{3}\)۔}

\medskip
\textbf{(ب)} \(|u|\le a\) اندر کی طرف بند ہوتا ہے، یعنی \(-a\le u\le a\):
\[
-8\le z^{2}-8\le 8 \;\Longrightarrow\; 0\le z^{2}\le 16
\]
بایاں حصہ \(z^{2}\ge0\) ہر حقیقی \(z\) کے لیے خود بخود درست ہے اور کوئی پابندی
نہیں لگاتا۔ دایاں حصہ دیتا ہے:
\[
z^{2}\le 16 \;\Longrightarrow\; -4\le z\le 4
\]

\jawab{(الف) \(y=-9,\ -\tfrac13\) \quad (ب) \(-4\le z\le 4\)}
\end{hal}

\begin{ghalti}
\(|u|\le a\) \emph{ایک} وقفہ دیتا ہے؛ \(|u|\ge a\) باہر کی طرف کھل کر
\emph{دو} وقفے دیتا ہے۔ ان دونوں کو الٹ دینے سے جواب سراسر الٹ جاتا ہے۔
\end{ghalti}

%---------------------------------------------------------------------
\section*{یونٹ \textenglish{4}: خطی مساوات}
\markboth{یونٹ 4: خطی مساوات}{}

\begin{nukta}[پہلے میلان، پھر ایک نقطہ]
\[
m=\frac{y_2-y_1}{x_2-x_1},\qquad
y-y_1=m(x-x_1),\qquad
y=mx+c
\]
متوازی خطوط: \(m_1=m_2\)۔ عمود خطوط: \(m_1m_2=-1\)۔
اس یونٹ کا ہر ''خط کی مساوات نکالیے'' والا سوال بس یہی دو قدم ہے۔
\end{nukta}

\begin{sawal}{سوال \textenglish{10}}
صابن کی طلب \(p=15-\tfrac76 q\) اور رسد \(p=\tfrac23 q\) سے دی گئی ہے، جہاں
\(p\) قیمت اور \(q\) مقدار ہے۔ توازنی قیمت اور توازنی مقدار نکالیے۔
\end{sawal}

\begin{hal}
منڈی وہاں صاف ہوتی ہے جہاں طلب اور رسد برابر ہو جائیں، یعنی جہاں دونوں خط ایک
دوسرے کو کاٹیں:
\[
15-\frac76 q=\frac23 q
\]
کسور ختم کرنے کے لیے پورے کو \(6\) سے ضرب دیجیے:
\[
90-7q=4q \;\Longrightarrow\; 90=11q \;\Longrightarrow\;
q^{*}=\frac{90}{11}\approx 8.18 \ \text{اکائیاں}
\]
\[
p^{*}=\frac23\cdot\frac{90}{11}=\frac{60}{11}\approx 5.45 \ \text{ڈالر}
\]
\parakh{طلب کی طرف سے:
\(15-\tfrac76\cdot\tfrac{90}{11}=15-\tfrac{105}{11}=\tfrac{60}{11}\)۔ دونوں
منحنی متفق ہیں، اس لیے تقاطع درست ہے۔}

\jawab{\(q^{*}=\tfrac{90}{11}\approx 8.18\) اکائیاں،
\(p^{*}=\tfrac{60}{11}\approx \$5.45\)}
\end{hal}

\begin{sawal}{سوال \textenglish{11}}
درج ذیل خصوصیات والے خط کی مساوات نکالیے:
\textbf{(الف)} \((-1,3)\) سے گزرتا ہو اور \(y=4x-5\) کے متوازی ہو؛
\textbf{(ب)} \((-5,4)\) سے گزرتا ہو اور \(2y=x+1\) پر عمود ہو؛
\textbf{(ج)} \((2,-8)\) سے گزرتا ہو اور \(x=-4\) کے متوازی ہو۔
\end{sawal}

\begin{hal}
\textbf{(الف)} متوازی خطوط کا میلان ایک ہوتا ہے، اس لیے \(m=4\):
\[
y-3=4(x+1) \;\Longrightarrow\; y=4x+7
\]

\textbf{(ب)} \(2y=x+1\) کو \(y=\tfrac12 x+\tfrac12\) لکھیے، اس کا میلان
\(\tfrac12\) ہے۔ عمود کا میلان منفی معکوس ہے:
\[
m=-\frac{1}{1/2}=-2
\;\Longrightarrow\;
y-4=-2(x+5) \;\Longrightarrow\; y=-2x-6
\]

\textbf{(ج)} \(x=-4\) ایک \emph{عمودی} خط ہے جس کا کوئی میلان ہوتا ہی نہیں۔ اس
کے متوازی خط بھی عمودی ہو گا اور دیے گئے نقطے کی \(x\) قدر سے گزرے گا:
\[
x=2
\]

\jawab{(الف) \(y=4x+7\) \quad (ب) \(y=-2x-6\) \quad (ج) \(x=2\)}
\end{hal}

\begin{ghalti}
حصہ (ج) اکثر امیدواروں کو پکڑ لیتا ہے۔ \(m_1m_2=-1\) کا اصول عمودی خطوط پر
لاگو نہیں ہوتا، کیونکہ عمودی خط کا میلان غیر معین ہے۔ جواب \(x=\)~ثابت ہوتا ہے،
اور وہ ثابت نقطے کی \(x\) قدر سے آتا ہے، \(y\) سے نہیں۔
\end{ghalti}

\begin{sawal}{سوال \textenglish{12}}
ایک کمپنی کرسیاں بناتی ہے۔ ہفتہ وار ثابت لاگت \textenglish{\$1{,}100} اور فی
کرسی متغیر لاگت \textenglish{\$40} ہے۔ \(x\) کرسیوں کی کل ہفتہ وار لاگت لکھیے۔
\textenglish{\$4{,}500} میں کتنی کرسیاں بن سکتی ہیں؟
\end{sawal}

\begin{hal}
خطی لاگت کا ڈھانچہ یہ ہے کہ متغیر لاگت \emph{میلان} بنتی ہے اور ثابت لاگت
\emph{قاطع}:
\[
C(x)=40x+1100
\]
\[
40x+1100=4500 \;\Longrightarrow\; 40x=3400 \;\Longrightarrow\; x=85
\]
\parakh{\(40(85)+1100=3400+1100=4500\)۔}

\jawab{\(C(x)=40x+1100\)؛ \textenglish{\$4{,}500} میں \textbf{\textenglish{85}
کرسیاں} بنتی ہیں۔}
\end{hal}

\begin{sawal}{سوال \textenglish{13}}
جان نے نئی گاڑی خریدی۔ \textenglish{7.5\%} سیلز ٹیکس اور \textenglish{\$150}
رجسٹریشن فیس ملا کر کل \textenglish{\$25{,}868.50} ادا کیے۔ گاڑی کی اصل قیمت
کیا تھی؟
\end{sawal}

\begin{hal}
فرض کیجیے ٹیکس سے پہلے قیمت \(P\) ہے۔ ٹیکس \(P\) پر لگتا ہے اور فیس بعد میں
جمع ہوتی ہے:
\[
P+0.075P+150=25{,}868.50
\]
\[
1.075P=25{,}718.50 \;\Longrightarrow\;
P=\frac{25{,}718.50}{1.075}=23{,}924.19
\]
\parakh{\(23{,}924.19+1{,}794.31+150=25{,}868.50\)۔}

\jawab{گاڑی کی اصل قیمت تقریباً \textbf{\textenglish{\$23{,}924.19}} تھی۔}
\end{hal}

\begin{ghalti}
کل رقم میں سے \textenglish{7.5\%} گھٹانا غلط ہے۔ \(25{,}718.50\times 0.925\)
سے \(23{,}789.61\) آتا ہے جو درست نہیں، کیونکہ ٹیکس بڑی (کل) رقم پر نہیں بلکہ
چھوٹی (اصل) قیمت پر لگا تھا۔ ہمیشہ \(1+r\) سے \emph{تقسیم} کیجیے،
\(1-r\) سے ضرب نہیں۔
\end{ghalti}

\clearpage
